% ==========================================================================
% Section 7: Discussion, Limitations, and Ethical Considerations
% ==========================================================================
\section{Discussion, Limitations, and Ethical Considerations}
\label{sec:discussion-limitations-ethics}

\subsection{Interpretation of the Main Findings}
\label{sec:interpretation-main-findings}

The ablations yield one clear positive retrieval intervention and several
conditional or negative findings. Cross-encoder reranking improves the rank of
evidence already admitted to the shared C30 pool; Jina has the strongest
retrieval result, while BGE has the strongest lexical answer overlap (the
paired evidence is linked in
\hyperref[tab:reranker-paired]{Table~\ref*{tab:reranker-paired}}). By
contrast, the tested hybrid and graph configurations reduce exact gold-chunk
coverage, as shown by the dense/hybrid comparison in
\hyperref[tab:dense-hybrid-results]{Table~\ref*{tab:dense-hybrid-results}} and
the graph displacement counts in
\hyperref[tab:graph-effect]{Table~\ref*{tab:graph-effect}}. These statements
concern a finite benchmark and its gold evidence, not legal advice quality in
deployment.

\subsection{Retrieval Quality, Context Coverage, and Latency Trade-offs}
\label{sec:retrieval-quality-tradeoffs}

DenseOnly is the preferred retrieval baseline in the dense/hybrid study: it
achieves Recall@10 of 0.4855 at 0.797-second retrieval p50. HybridOnly falls
to 0.4257 and takes 27.855 seconds when amortized BM25 time is included
(\hyperref[tab:dense-hybrid-results]{Table~\ref*{tab:dense-hybrid-results}}).
Graph contexts appear more identity-precise because they contain only
5.605--6.305 chunks on average, but their reduced context size also removes
gold coverage (\hyperref[tab:graph-effect]{Table~\ref*{tab:graph-effect}}).
The reranker result has a different interpretation: it operates after a fixed
candidate ceiling, so Jina's gain is a ranking gain rather than proof that the
underlying candidate retriever finds more evidence.

\subsection{Implications for Legal RAG Design}
\label{sec:legal-rag-implications}

For this benchmark, use header-plus-chunk embeddings when correct document
identification is more valuable than the rank of an exact chunk; otherwise
prefer the chunk-only representation
(\hyperref[tab:embedding-results]{Table~\ref*{tab:embedding-results}}).
Do not deploy the current BM25+RRF or one-hop graph setup as the default
retrieval path without retuning. Prefer Jina when retrieval quality justifies
its 1.278-second reranking p50, or mMiniLM when its lower 0.385-second cost is
more important. Retain the Base prompt for lexical answer overlap. In all
cases, present citations as sources to inspect and retain human review for
legal interpretation.

\subsection{Dataset and Coverage Limitations}
\label{sec:dataset-coverage-limitations}

The benchmark has 500 questions, of which only 400 are answerable by stored
gold chunks. The hard stratum has one answerable item, while multi-hop has 20;
those slices are descriptive and cannot establish stable subgroup effects.
Gold chunk IDs are an evidence target, not necessarily the full set of legally
acceptable authorities. Corpus coverage, language variation, temporal validity,
and external legal references may all limit transfer beyond the data snapshot.

\subsection{Knowledge Graph and Validity Limitations}
\label{sec:graph-validity-limitations}

The graph study evaluates a specific policy--five seeds, one hop, and ten
retained chunks--rather than graph information in general. Its failures point
to an engineering hypothesis: expansion substitutes for seeds too early.
Future graph experiments should preserve high-ranked seeds, separately budget
neighbors, verify relationship direction and validity semantics, and measure
whether graph-only evidence is legally useful rather than merely novel.

\subsection{Retrieval and Generation Limitations}
\label{sec:retrieval-generation-limitations}

The index, BM25 tokenization, RRF parameter, candidate budgets, rerankers,
generator, prompt, and final context size are fixed settings, not optimized
universally. Reranker timings use cached BM25 candidates and omit full-online
sparse retrieval. Lexical overlap may reward reference phrasing rather than
legal adequacy. Structural citations measure marker presence or parser coverage,
not the relationship between a claim and the cited authority.

\subsection{Evaluation Limitations}
\label{sec:evaluation-limitations}

No blinded legal-correctness labels, claim-level faithfulness labels, or
validated semantic judge are available. The report therefore does not claim
that any configuration is legally more correct, safer, or more faithful solely
because it improves F1, ROUGE-L, citation count, or refusal markers. Full-online
p95 reranker latency, peak memory, and throughput are also absent. Non-GPT
prompt pairs have provenance or control differences and remain descriptive.

\subsection{Ethical Considerations and Safe Use}
\label{sec:ethical-safe-use}

L\_RAG should support research and source discovery, not replace a qualified
legal professional. Users should inspect the retrieved authority, its date and
validity, and the uncertainty of an answer. Systems should avoid unnecessary
personal query data, should not obscure abstentions, and should treat
well-cited generated text as a review cue rather than a legal conclusion.
